\section{From paths to monochromatic copies}\label{paths:section}

We now give a criterion that converts products of scaled copies into the
Ramsey property. The products will first be taken in a free group. A
word construction then turns their algebraic relations into actual
monochromatic point sets. Throughout this section,
$A=\{a_1,\ldots,a_s\}$ is a finite Euclidean configuration with $s\ge2$.

Let $\mathcal E=\R^{(\N)}$ be the Euclidean space of real sequences with
finite support, and let $K$ be the free group on the symbols
$\xi_z$, indexed by $z\in\mathcal E$. Write
$\Delta g=(g,\ldots,g)\in K^s$ for $g\in K$.
A \emph{path} is a chosen finite sequence of factors in $K^s$ of the
following two kinds:
\begin{enumerate}
\item a diagonal factor $\Delta g$, assigned weight zero;
\item a factor $X_b=(\xi_{b_1},\ldots,\xi_{b_s})$, where
      $b_1,\ldots,b_s\in\mathcal E$ and, for some $u\ge0$,
      \[
      \|b_i-b_j\|^2=u\|a_i-a_j\|^2\qquad(1\le i,j\le s),
      \]
      assigned weight $u$.
\end{enumerate}
The \emph{endpoint} is the product of the factors in $K^s$, in their
given order. The \emph{weight} is the sum of their assigned weights.
The weight belongs to the chosen path, not to its endpoint. In
particular, two paths with the same endpoint can have different weights.

\begin{proposition}[A path criterion]\label{paths:criterion}
Suppose that for every integer $n\ge1$ there are two paths with the
same endpoint and weights $\lambda>0$ and $\mu\ge0$ satisfying
$\mu/\lambda<1/n$. Then $A$ is Euclidean Ramsey.
\end{proposition}

We prove the proposition in three steps. We first reduce to colorings
of $\mathcal E$ and synchronize the available path weights. We then
develop the ultrafilter facts needed to realize group products by
finite strings. Finally, the Hales--Jewett theorem chooses the variable
positions of a monochromatic string family.

\subsection{Finite dimensions and synchronized weights}

The following finite-dimensional reduction is the compactness principle
used in Euclidean Ramsey theory; compare
\cite[Proposition~4]{egm-rss}. We include its proof for arbitrary colorings.

\begin{lemma}\label{paths:compactness}
Fix an integer $r\ge2$. If every $r$-coloring of $\mathcal E$ contains
a monochromatic congruent copy of $A$, then the same holds for every
$r$-coloring of $\R^D$, for some finite $D$.
\end{lemma}

\begin{proof}
Suppose instead that every $\R^D$ admits an $r$-coloring avoiding $A$.
Consider the compact product space $\{1,\ldots,r\}^{\mathcal E}$.
For each ordered congruent copy $(b_1,\ldots,b_s)$ of $A$ in
$\mathcal E$, impose the condition that the colors of the $b_i$ are
not all equal. This is a clopen condition depending on finitely many
coordinates of the product space.

Every finite collection of these conditions involves finitely many
finite-support vectors. They all lie in the first $D$ coordinate
positions for some $D$. An avoiding coloring of this $\R^D$, extended
arbitrarily to $\mathcal E$, satisfies the finite collection.
Compactness therefore gives a coloring satisfying all the conditions,
contrary to the hypothesis.
\end{proof}

\begin{lemma}\label{paths:synchronization}
Fix an integer $n\ge1$. Let two paths have the same endpoint and weights $\lambda>0$ and
$\mu\ge0$. If $\rho=\mu/\lambda<1/n$, then there is an endpoint
$w\in K^s$ admitting a path of weight $1/t$ for every integer
$1\le t\le n$.
\end{lemma}

\begin{proof}
Set
\[
\alpha_t=\frac{1/t-\rho}{1-\rho}\qquad(1\le t\le n).
\]
These numbers satisfy
$0<\alpha_n<\cdots<\alpha_1=1$. Partition $[0,1]$ at these points.
For each resulting interval of length $l>0$, apply to both paths the
endomorphism of $K$ defined by
\[
\xi_z\longmapsto\xi_{\sqrt{l/\lambda}\,z}.
\]
It preserves diagonal factors and multiplies every copy-factor weight
by $l/\lambda$. The two new paths thus have the same endpoint and
weights $l$ and $\rho l$.

Concatenate the interval paths in increasing interval order. The
resulting endpoint is independent of which of the two paths is chosen
on each interval. For a given $t$, choose the larger weight on the
intervals before $\alpha_t$ and the smaller weight on the remaining
intervals. The total weight is
\[
\alpha_t+\rho(1-\alpha_t)=\frac1t.
\]
The fixed order of concatenation gives one endpoint $w$ for all $t$;
no commutation of factors is used.
\end{proof}

\subsection{Realizing ultrafilter products by strings}

Let $\mathcal S$ be the monoid of finite strings on the alphabet
$\mathcal E$, with concatenation as multiplication and the empty
string $\epsilon$ as identity. Choose pairwise disjoint infinite
coordinate blocks in $\mathcal E$, and let
$J_j:\mathcal E\to\mathcal E$ be the corresponding linear isometric
embeddings, one for each $j\ge1$. For a finite string
$z=(z_1,\ldots,z_L)$, define
\[
\Phi(z)=\sum_{j=1}^L J_jz_j.
\]
For two strings of the same length,
\begin{equation}\label{paths:string-distance}
\|\Phi(z)-\Phi(z')\|^2
=\sum_{j=1}^L\|z_j-z'_j\|^2.
\end{equation}
Thus we seek a realization in which each copy factor supplies its
prescribed one-letter strings, while each diagonal factor supplies
one common string in all $s$ roles. Common strings will contribute
zero to the squared distances; the remaining contribution will be
the path weight times the corresponding squared distance in $A$.
Diagonal factors may contain group inverses, so they cannot in general
be evaluated directly in the string monoid. We instead map $K$ into a
group of ultrafilters on $\mathcal S$ and then realize its common
factors by common strings.

The following standard construction fixes the multiplication convention; see
\cite[Theorem~2.5, Corollary~2.6, and Theorem~1.42]{hindman-strauss}.
The classical idempotent argument also appears in Ellis's semigroup lemma
\cite[Lemma~1]{ellis-idempotent}; taking the opposite semigroup matches
the one-sided continuity convention used here.

An ultrafilter on $\mathcal S$ is a maximal proper filter of subsets
of $\mathcal S$. Write $\beta\mathcal S$ for the space of all such
ultrafilters. For $B\subseteq\mathcal S$, its basic clopen set is
\[
\overline B=\{u\in\beta\mathcal S:B\in u\}.
\]
This space is compact and Hausdorff: the ultrafilter conditions define a closed
subset of $\{0,1\}^{\mathcal P(\mathcal S)}$.
For a string $x$, put $x^{-1}B=\{y\in\mathcal S:xy\in B\}$;
this notation denotes a preimage under concatenation, not a string
inverse. Define the product of ultrafilters by
\begin{equation}\label{paths:ultrafilter-product}
B\in uv
\quad\Longleftrightarrow\quad
\{x\in\mathcal S:x^{-1}B\in v\}\in u.
\end{equation}
The principal ultrafilter at $x$ is denoted by $\delta_x$.

\begin{lemma}\label{paths:corner-group}
The operation in \eqref{paths:ultrafilter-product} makes
$\beta\mathcal S$ a compact right-topological semigroup with identity
$\delta_\epsilon$. It contains an idempotent $p$ such that
$p(\beta\mathcal S)p$ is a group with identity $p$.
\end{lemma}

\begin{proof}
For fixed $v$, the map $B\mapsto\{x:x^{-1}B\in v\}$ preserves
finite intersections and complements. Thus
\eqref{paths:ultrafilter-product} defines an ultrafilter $uv$.
For $u,v,w\in\beta\mathcal S$, both parenthesizations of their
product contain $B$ precisely when
\[
\{x:\{y:\{z:xyz\in B\}\in w\}\in v\}\in u.
\]
This proves associativity. For fixed $v$, the inverse image of
$\overline B$ under $u\mapsto uv$ is
$\overline{\{x:x^{-1}B\in v\}}$, so right translation is
continuous. The assertion about $\delta_\epsilon$ follows directly
from the definition.

Choose a minimal nonempty compact left ideal $L_0$ in
$\beta\mathcal S$. Such an ideal exists because every chain of
nonempty compact left ideals has a nonempty intersection. For each
$v\in L_0$, the set $(\beta\mathcal S)v$ is a nonempty compact
left ideal contained in $L_0$, hence
\begin{equation}\label{paths:minimal-left-ideal}
(\beta\mathcal S)v=L_0.
\end{equation}
The ideal $L_0$ is also a compact subsemigroup.

Choose a minimal nonempty compact subsemigroup $H_0\subseteq L_0$.
For $x\in H_0$, the set $H_0x$ is compact and is a subsemigroup:
$(ax)(bx)=a(xb)x\in H_0x$ for $a,b\in H_0$.
Minimality gives $H_0x=H_0$. Consequently
\[
\{y\in H_0:yx=x\}
\]
is nonempty, compact, and a subsemigroup. It equals $H_0$ by
minimality, so $x^2=x$. In particular $L_0$ contains an idempotent
$p$.

The corner $G=p(\beta\mathcal S)p$ is a monoid with identity $p$.
Every $y\in G$ lies in $L_0$, so \eqref{paths:minimal-left-ideal}
gives $p=vy$ for some $v\in\beta\mathcal S$. Since $py=y$,
\[
(pvp)y=pvy=p.
\]
Thus every element of $G$ has a left inverse in $G$. A monoid with
this property is a group: if $ay=p$ and $ba=p$, then
$b=b(ay)=(ba)y=y$, whence $ya=p$ as well.
\end{proof}

The next lemma explains why common ultrafilter factors can be replaced
by common strings. This is the step that will preserve the distances
between the different roles in a monochromatic configuration.

\begin{lemma}[Simultaneous realization]\label{paths:realization}
Let $q_{i,j}\in\beta\mathcal S$ for $1\le i\le s$ and
$1\le j\le M$, where $M\ge1$. Suppose that every column $j$ has
one of the following specified forms:
\begin{enumerate}
\item $q_{i,j}=q_j$ is independent of $i$;
\item $q_{i,j}=\delta_{z_{i,j}}$ for prescribed strings
      $z_{i,j}\in\mathcal S$.
\end{enumerate}
If $B\in q_{i,1}\cdots q_{i,M}$ for every $i$, there are strings
$x_{i,j}$ such that all products $x_{i,1}\cdots x_{i,M}$ belong to
$B$, with $x_{i,j}$ independent of $i$ in each column of the first
form and $x_{i,j}=z_{i,j}$ in each column of the second form.
\end{lemma}

\begin{proof}
Choose the strings from left to right. Before column $j$, let $P_i$
be the already chosen prefix and maintain the invariant
\[
P_i^{-1}B\in q_{i,j}\cdots q_{i,M}.
\]
Initially $P_i=\epsilon$. Put
$v_i=q_{i,j+1}\cdots q_{i,M}$, interpreting an empty product here
as $\delta_\epsilon$. By
\eqref{paths:ultrafilter-product}, the set
\[
E_i=\{x:x^{-1}(P_i^{-1}B)\in v_i\}
\]
belongs to $q_{i,j}$. In a common column, the finite intersection
$\bigcap_i E_i$ belongs to the common ultrafilter and is nonempty;
choose the same string from it for all $i$. In a prescribed principal
column, $z_{i,j}\in E_i$, so choose that string. Replacing $P_i$ by
$P_ix_{i,j}$ preserves the invariant. After the last column it says
$P_i^{-1}B\in\delta_\epsilon$, which is precisely $P_i\in B$.
\end{proof}

We now have both ingredients for the word construction: a group in
which the path endpoints can be evaluated, and a way to realize
their products by strings while choosing every common factor equally
in all roles.

\subsection{The Hales--Jewett step}

\begin{proof}[Proof of Proposition~\ref{paths:criterion}]
Fix $r\ge2$ and an arbitrary coloring
$c:\mathcal E\to\{1,\ldots,r\}$. Choose $n$ such that every
$r$-coloring of $\{1,\ldots,s\}^n$ has a monochromatic
combinatorial line, by the Hales--Jewett theorem
\cite{hales-jewett}; see also Shelah's proof and primitive-recursive
bound \cite[p.~686 and Section~1]{shelah}. We use only existence of
the length. Such a line has a nonempty set of variable
positions, occupied by one common letter, and fixed letters in all
other positions.

In the group formulation of Leader, Russell and Walters
\cite[Conjecture~C and the following discussion]{lrw-transitive},
the number of varying positions is prescribed before coloring;
without that restriction, ordinary Hales--Jewett applies.
Here we allow the line to choose that number and prepare all the
required path weights in advance.
The hypothesis and Lemma~\ref{paths:synchronization} give an endpoint
$w=(w_1,\ldots,w_s)\in K^s$ and a path of weight $1/t$ ending at
$w$ for every $1\le t\le n$.

Color $\mathcal S$ by $c\circ\Phi$, and write its color classes
as $B_1,\ldots,B_r$.

Choose $p$ as in Lemma~\ref{paths:corner-group}. The freeness of $K$
gives a homomorphism
\[
\psi:K\longrightarrow p(\beta\mathcal S)p,
\qquad \psi(\xi_z)=p\delta_{(z)}p,
\]
where $(z)$ is the one-letter string with entry $z$.
Color a word $i_1\cdots i_n$ by the unique index $k$ for which
\[
B_k\in\psi(w_{i_1})\cdots\psi(w_{i_n}).
\]
Every ultrafilter contains exactly one class of a finite partition,
so this defines an $r$-coloring.

Take a monochromatic combinatorial line and let $t$ be its number
of variable positions. Figure~\ref{fig:synchronization} records the
order of these choices.

\begin{figure}[tbp]
\centering
\begin{tikzpicture}[
  box/.style={draw=blue!45!black,rounded corners=2pt,fill=blue!3,
    text width=5.6cm,minimum height=1.2cm,align=center,
    inner sep=5pt,font=\small},
  flow/.style={-{Stealth[length=2mm]},thick,draw=blue!45!black}
]
\node[box] (length) {Fix the number of colors $r$.\\
  Choose a Hales--Jewett length $n$.};
\node[box,right=9mm of length] (endpoint) {Prepare one endpoint $w$\\
  with paths of weights $1,\frac12,\ldots,\frac1n$.};
\node[box,below=10mm of endpoint] (line) {Color words using $w$.\\
  A monochromatic line has $t$ variable positions.};
\node[box,left=9mm of line] (scale) {Use the path of weight $1/t$\\
  at each variable position.\\
  Squared-distance factor: $t(1/t)=1$.};
\draw[flow] (length) -- (endpoint);
\draw[flow] (endpoint) -- (line);
\draw[flow] (line) -- (scale);
\end{tikzpicture}
\caption{The order of choices in the path criterion. Every possible
variable count is accommodated before the word coloring is defined.
Choosing a path representation later leaves its endpoint unchanged;
common factors become common strings, so only the variable blocks
contribute to pairwise distances.}
\label{fig:synchronization}
\end{figure}


For each of the line's $s$ words, expand the factor
at every variable position using the chosen path ending at $w$ with
weight $1/t$. A diagonal path factor contributes one common
ultrafilter $\psi(g)$. A copy factor contributes
$p\delta_{(b_i)}p$ in role $i$: the two $p$ factors are common and
the middle factor is principal at a prescribed one-letter string.
Unexpanded fixed positions are common factors as well. The identity
of the corner group is $p$; if an empty group product occurs, it is
retained as this common factor. The expanded products therefore have
the column forms required by Lemma~\ref{paths:realization}, and all
contain the same color class $B_k$.

Apply that lemma. It produces strings $z^{(1)},\ldots,z^{(s)}$ in $B_k$. Every common
factor was realized by the same string in all roles, and every
varying factor has length one. Hence the resulting strings have equal
lengths and their entries align. Common strings contribute zero to
their pairwise squared distances. Each of the $t$ expanded paths
contributes
$\|a_i-a_j\|^2/t$ by its weight. Equation
\eqref{paths:string-distance} consequently gives
\[
\|\Phi(z^{(i)})-\Phi(z^{(j)})\|^2
=t\frac1t\|a_i-a_j\|^2
=\|a_i-a_j\|^2.
\]
Their images form a monochromatic congruent copy of $A$ in
$\mathcal E$. Lemma~\ref{paths:compactness} now supplies a finite
forcing dimension for this $r$. Since $r$ was arbitrary, $A$ is
Euclidean Ramsey.
\end{proof}
